\documentclass[11pt]{amsart}

\usepackage{amsmath}
\usepackage{amssymb}
\usepackage{amsthm}
%\usepackage{psfig}

\begin{document}
\baselineskip=16pt
\textheight=8.5in
%\parindent=0pt 
\def\sk {\hskip .5cm}
\def\skv {\vskip .05cm}
\def\cos {\mbox{cos}}
\def\sin {\mbox{sin}}
\def\tan {\mbox{tan}}
\def\intl{\int\limits}
\def\lm{\lim\limits}
\newcommand{\frc}{\displaystyle\frac}
\def\xbf{{\mathbf x}}
\def\fbf{{\mathbf f}}
\def\gbf{{\mathbf g}}

\def\dbA{{\mathbb A}}
\def\dbB{{\mathbb B}}
\def\dbC{{\mathbb C}}
\def\dbD{{\mathbb D}}
\def\dbE{{\mathbb E}}
\def\dbF{{\mathbb F}}
\def\dbG{{\mathbb G}}
\def\dbH{{\mathbb H}}
\def\dbI{{\mathbb I}}
\def\dbJ{{\mathbb J}}
\def\dbK{{\mathbb K}}
\def\dbL{{\mathbb L}}
\def\dbM{{\mathbb M}}
\def\dbN{{\mathbb N}}
\def\dbO{{\mathbb O}}
\def\dbP{{\mathbb P}}
\def\dbQ{{\mathbb Q}}
\def\dbR{{\mathbb R}}
\def\dbS{{\mathbb S}}
\def\dbT{{\mathbb T}}
\def\dbU{{\mathbb U}}
\def\dbV{{\mathbb V}}
\def\dbW{{\mathbb W}}
\def\dbX{{\mathbb X}}
\def\dbY{{\mathbb Y}}
\def\dbZ{{\mathbb Z}}

\def\la{{\langle}}
\def\ra{{\rangle}}
\def\summ{{\sum\limits}}
\def\eps{{\varepsilon}}
\def\lam{{\lambda}}
\def\tr{{\rm tr}}
\def\char{{\rm char}}
\def\Mat{{\rm Mat}}
\def\Func{{\rm Func}}
\def\Span{{\rm Span}}
\def\Ker{{\rm Ker}}
\def\Im{{\rm Im}}



\bf\centerline{Homework \#4. Due on Thursday, Sep 24th, 11:59pm on Canvas}\rm
\vskip .1cm

\bf{Reading for this homework: }\rm Sections 2.1, 2.2 and parts of 2.4.

\bf{Plan for next week (Sep 21,23): }\rm  2.3 (composition of linear maps and matrix multiplication), 2.5 (the change of coordinate matrix),
parts of 2.4 (invertibility and isomorphisms). Note that we have already proved some results from 2.4 in Lecture~7 (Sep 16).
\skv

\bf\centerline{Problems: }\rm
\skv
\bf{Problem 1: }\rm This problem illustrates a relation between the internal
and external direct sums (see also Problem~3 below). Let $U$ and $W$ be vector spaces over the same field and let $V=U\oplus W$ be their external direct sum, that is, $V$ is the set
of ordered pairs $\{(u,w): u\in U, w\in W\}$ with componentwise addition and
scalar multiplication: $(u_1,w_1)+(u_2,w_2)=(u_1+u_2,w_1+w_2)$ for all $u_1,u_2\in U$ and $w_1,w_2\in W$
and $\lam (u,w)=(\lam u,\lam w)$ for all $\lam\in F$, $u\in U$ and 
$w\in W$.
 Define the subsets $\widetilde U$ and $\widetilde W$ of $V$
by $\widetilde U=\{(u,0) : u\in U\}$ and $\widetilde W=\{(0,w) : u\in U\}$.
Prove that
\begin{itemize}
\item[(i)] $\widetilde U$ and $\widetilde W$ are subspaces of $V$
isomorphic to $U$ and $W$, respectively;
\item[(ii)] $V=\widetilde U\oplus \widetilde W$, where now we talk
about the internal direct sum.
\end{itemize}
Conclude that the external direct sum of two spaces $U$ and $W$ is equal to
the internal direct sum of isomorphic copies of $U$ and $W$.
\skv

\noindent
\bf{Problem 2: }\rm Let $V$ be a vector space over $F$,
and let $V^2=V\oplus V$ be the external direct sum of two copies of $V$.
Fix $\lam\in F$ and put $W=\{(v,\lam v) : v\in V\}$. Prove that 
$W$ is a subspace of $V^2$ which is isomorphic to $V$. 

\skv
\bf{Problem 3: }\rm Let $U$ and $W$ be subspaces of a finite-dimensional 
vector space $V$.
Consider the external direct sum $U\oplus W=\{(u,w) : u\in U, w\in W\}$
and the map $T:U\oplus W\to V$ given by $T((u,w))=u+w$. 
\begin{itemize}
\item[(a)] Prove that $T$ is linear and use the rank-nullity theorem to prove that $$\dim(U+W)=\dim(U)+\dim(W)-\dim(U\cap W)$$ 
(recall that this was proved by a much longer argument in HW\#3.4).
\item[(b)] Show that if the internal sum $U+W$ is direct, then
$T$ is an isomorphism between the external and internal
direct sums of $U$ and $W$.
\end{itemize}
\skv
\bf{Problem 4: }\rm Recall that $\mathfrak{sl}_n(F)$ denotes
the space of all $n\times n$ matrices over $F$ with trace $0$.
In Problem~6 of HW\#2 it was proved that $\dim(\mathfrak{sl}_n(F))=n^2-1$
after a considerable amount of work. Now give a short proof of this
fact by applying the rank-nullity theorem to a suitable linear map.
\skv
\skv
{\bf Problem 5: }Let $V$ be a finite-dimensional vector space over a field $F$.
\begin{itemize}
\item[(a)] Let $T:V\to V$ be a linear map. Prove that $T$ is an injective $\iff$ $T$ is surjective.
\item[(b)] Now suppose that $c_0,\ldots, c_n$ are distinct elements of $F$ and consider the map
$T:\mathcal P_n(F)\to F^{n+1}$ (where as before $\mathcal P_n(F)$ is polynomials of degree $\leq n$)
given by $T(p(x))=(p(c_0),p(c_1),\ldots, p(c_n))$. Use (a) to prove that $T$ is an isomorphism.
You may use without proof that for a nonzero polynomial of degree $n$ cannot have more than $n$ roots.
\item[(c)] Give examples showing that both implications in (a) become false if we do not assume that $V$ is finite-dimensional.
\end{itemize}
\skv

\noindent
\bf{Problem 6: }\rm Let $V$ be a vector space and $T:V\to V$ a linear map.
A subspace $W$ of $V$ is called {\it $T$-invariant} if $T(W)\subseteq W$
where $T(W)=\{T(w): w\in W\}$.
\begin{itemize}
\item[(a)] Prove that $\Ker(T)$ and $\Im(T)$ are $T$-invariant subspaces.
\item[(b)] Assume that $\dim(V)<\infty$ and $W$ is a $T$-invariant subspace
of $V$ s.t. $V=W\oplus \Ker(T)$. Prove that $W=\Im(T)$. {\bf Hint:}
First show that $\Im(T)\subseteq W$.
\item[(c)] Give an example with $\dim(V)<\infty$ where the sum
$\Im(T)+\Ker(T)$ is NOT direct. 
\item[(d)] Use (b) and (c) to conclude that a $T$-invariant subspace
may NOT have a $T$-invariant complement.
\end{itemize}

\skv


\noindent
{\bf Problem 7: } Let $V=\mathcal P_2(\dbR)$, $\beta=\{1,x,x^2\}$ and $\gamma=\{1,x^2-x+1, x^2+x+1\}$.
\begin{itemize}
\item[(a)] Prove that $\gamma$ is a basis for $V$
\item[(b)] Now consider the linear map $T:V\to V$ given by $$T(p(x))=p(x)-p'(x).$$ Compute the matrices
$[T]_{\beta}^{\gamma}$ and $[T]_{\gamma}^{\beta}$ directly from definition. Include all the details.
\end{itemize}
\skv
{\bf Problem 8: } Let $V$ be a finite-dimensional vector space, $W$ a subspace of $V$,
$n=\dim(V)$ and $k=\dim(W)$. Choose a basis $\beta_W=\{w_1,\ldots, w_k\}$ of $W$
and extend it to a basis $\beta=\{w_1\ldots, w_k, u_1,\ldots, u_{n-k}\}$ of $V$.

Find (with proof) a simple condition (***) on the entries of an $n\times n$ matrix $A$
over $F$ such that the following statement holds:

$W$ is $T$-invariant (as defined in Problem~6) $\iff$ the matrix $[T]_{\beta}=[T]_{\beta}^{\beta}$
satisfies (***).
\end{document}